\section*{Chapter \ref{ch:g11:organic-skeletons} --- Organic Molecules: Skeletons and Names}

\begin{solution}{exo:g11:organic-skeletons:1}
Both have six carbon \omterm{def:g7:atoms-and-molecules:atom}{atoms} (ends and corners) and \ce{C6H14}: hexane and
2-methylpentane are \omterm{def:g11:organic-skeletons:isomer}{isomers}.
\end{solution}

\begin{solution}{exo:g11:organic-skeletons:2}
Propane; hexane; 2-methylbutane.
\end{solution}

\begin{solution}{exo:g11:organic-skeletons:3}
\ce{CH3-CH2-CH3}; \chemfig{CH_3-CH(-[2]CH_3)-CH_3};
\chemfig{H-C(-[2]H)(-[6]H)-C(-[2]H)(-[6]H)-H}.
\end{solution}

\begin{solution}{exo:g11:organic-skeletons:4}
The first, 3-methylpentane: 6 carbon \omterm{def:g7:atoms-and-molecules:atom}{atoms}; its \ce{CH3} ends (three of
them) carry 9 hydrogen \omterm{def:g7:atoms-and-molecules:atom}{atoms}, its two \ce{CH2} corners 4 and its
\ce{CH} branch point 1: \ce{C6H14}. Cyclohexane: 6 corners, each
\ce{CH2}: \ce{C6H12}.
\end{solution}

\begin{solution}{exo:g11:organic-skeletons:5}
\ce{C8H18}. $2n + 2 = 20$ gives $n = 9$: nonane, \ce{C9H20}.
\end{solution}

\begin{solution}{exo:g11:organic-skeletons:6}
\setchemfig{atom sep=1.6em}
3-methylhexane \chemfig{-[:30]-[:-30]-[:30](-[:90])-[:-30]-[:30]-[:-30]},
\ce{C7H16}; 2,2-dimethylbutane \chemfig{-[:0](-[:90])(-[:-90])-[:0]-[:30]},
\ce{C6H14}; 3-ethylpentane
\chemfig{-[:30]-[:-30](-[:-90]-[:-30])-[:30]-[:-30]}, \ce{C7H16}.
\end{solution}

\begin{solution}{exo:g11:organic-skeletons:7}
Pentane, 2-methylbutane and 2,2-dimethylpropane, drawn in the figure of
the chapter.
\end{solution}

\begin{solution}{exo:g11:organic-skeletons:8}
``2-ethylbutane'' has a longest chain of five carbons through the ethyl
group: 3-methylpentane. ``4-methylpentane'' is numbered from the wrong
end: 2-methylpentane. ``1-methylbutane'' is simply a chain of five
carbons: pentane.
\end{solution}

\begin{solution}{exo:g11:organic-skeletons:9}
2,2-dimethylpropane (\qty{9.5}{\celsius}) < 2-methylbutane
(\qty{27.8}{\celsius}) < pentane (\qty{36.1}{\celsius}). The same \omterm{def:g7:atoms-and-molecules:atom}{atoms}
are arranged more and more compactly: less contact between \omterm{def:g7:atoms-and-molecules:molecule}{molecules},
weaker \omterm{def:g11:polarity-and-cohesion:van-der-waals}{van der Waals interactions}, lower boiling temperature.
\end{solution}

\begin{solution}{exo:g11:organic-skeletons:10}
Closing the chain into a ring makes one more \ce{C-C} bond, which takes
the place of two \ce{C-H} bonds (one at each end). Different \omterm{def:g11:organic-skeletons:formulas}{molecular
formulas}: not \omterm{def:g11:organic-skeletons:isomer}{isomers}. Cyclohexane is not an \omterm{def:g11:organic-skeletons:alkane}{alkane} (its skeleton is a
ring; its formula is not \ce{C_{n}H_{2n+2}}).
\end{solution}

\begin{solution}{exo:g11:organic-skeletons:11}
Butane and 2-methylpropane (\ce{C4H10}). Propane is \ce{C3H8} and
cyclobutane \ce{C4H8}.
\end{solution}

\begin{solution}{exo:g11:organic-skeletons:12}
\setchemfig{atom sep=1.5em}
\begin{center}
\small
\begin{tabular}{ccc}
\chemfig{-[:30]-[:-30]-[:30]-[:-30]-[:30]} &
\chemfig{-[:30](-[:90])-[:-30]-[:30]-[:-30]} &
\chemfig{-[:30]-[:-30](-[:-90])-[:30]-[:-30]} \\
hexane & 2-methylpentane & 3-methylpentane \\[6pt]
\chemfig{-[:0](-[:90])(-[:-90])-[:0]-[:30]} &
\chemfig{-[:30](-[:90])-[:-30](-[:-90])-[:30]} & \\
2,2-dimethylbutane & 2,3-dimethylbutane &
\end{tabular}
\end{center}
\end{solution}

\begin{solution}{exo:g11:organic-skeletons:13}
Two long zigzags can lie side by side along their whole length, with
many \omterm{def:g7:atoms-and-molecules:atom}{atoms} in contact; two balls touch at only a small area. \omterm{def:g11:polarity-and-cohesion:van-der-waals}{Van der
Waals interactions} act between \omterm{def:g7:atoms-and-molecules:atom}{atoms} in contact: they are larger
between linear \omterm{def:g7:atoms-and-molecules:molecule}{molecules}, which need a higher temperature to separate.
\end{solution}

\begin{solution}{exo:g11:organic-skeletons:14}
The longest chain has six carbons; numbered from the left the
substituents are at 2, 4, 4, from the right at 3, 3, 5; the first
difference (2 against 3) decides: 2,4,4-trimethylhexane.
\end{solution}

\begin{solution}{exo:g11:organic-skeletons:15}
\setchemfig{atom sep=1.6em}
\chemfig{-[:0](-[:90])(-[:-90])-[:-30]-[:30](-[:90])-[:-30]}: a chain of
five carbons with two methyl groups on carbon 2 and one on carbon 4,
eight carbon \omterm{def:g7:atoms-and-molecules:atom}{atoms} in all; $2 \times 8 + 2 = 18$ hydrogen \omterm{def:g7:atoms-and-molecules:atom}{atoms},
\ce{C8H18}, the formula of octane: they are \omterm{def:g11:organic-skeletons:isomer}{isomers}. \omterm{def:g11:organic-skeletons:formulas}{Semi-structural
formula}: \chemfig{CH_3-C(-[2]CH_3)(-[6]CH_3)-CH_2-CH(-[2]CH_3)-CH_3}.
\end{solution}

\begin{solution}{pb:g11:organic-skeletons:1}
\setchemfig{atom sep=1.5em}
\textbf{1.} \ce{C4H10}: with $n = 4$, $2n + 2 = 10$.

\textbf{2.} \chemfig{H-C(-[2]H)(-[6]H)-C(-[2]H)(-[6]H)-C(-[2]H)(-[6]H)-C(-[2]H)(-[6]H)-H}.

\textbf{3.} \ce{CH3-CH2-CH2-CH3}; \chemfig{-[:30]-[:-30]-[:30]}.

\textbf{4.} \chemfig{H-C(-[2]H)(-[6]H)-C(-[6]H)(-[2,1.5]C(-[1]H)(-[3]H)(-[2]H))-C(-[2]H)(-[6]H)-H};
\chemfig{CH_3-CH(-[2]CH_3)-CH_3}; \chemfig{-[:30](-[:90])-[:-30]}. It
has the same \omterm{def:g7:atoms-and-molecules:atom}{atoms}, four carbons and ten hydrogens, bonded in a
different order.

\textbf{5.} Yes: \omterm{def:g11:organic-skeletons:isomer}{constitutional isomers} (different skeletons).

\textbf{6.} All three: their boiling temperatures are below
\qty{20}{\celsius}.

\textbf{7.} At \qty{-5}{\celsius}, below its boiling temperature, butane
stays liquid at normal pressure and gives off little gas.

\textbf{8.} Propane (\qty{-42}{\celsius}) and 2-methylpropane
(\qty{-11.7}{\celsius}) boil below winter temperatures: they still give
plenty of gas in the cold.

\textbf{9.} 2-methylpropane is branched, more compact: weaker \omterm{def:g11:polarity-and-cohesion:van-der-waals}{van der
Waals interactions}.

\textbf{10.} With three carbon \omterm{def:g7:atoms-and-molecules:atom}{atoms} the only possible skeleton is the
chain \ce{C-C-C}: any carbon moved elsewhere gives the same chain.

\textbf{11.} 2-methylbutane.

\textbf{12.} The chain must be numbered from the end nearer the
substituent: 2, not 3.

\textbf{13.} 2, 3 and 5.

\textbf{14.} Heptane, \chemfig{-[:30]-[:-30]-[:30]-[:-30]-[:30]-[:-30]}.

\textbf{15.} 2-methylhexane \chemfig{-[:30](-[:90])-[:-30]-[:30]-[:-30]-[:30]}
and 3-methylhexane \chemfig{-[:30]-[:-30](-[:-90])-[:30]-[:-30]-[:30]}.
``4-methylhexane'' is 3-methylhexane numbered from the wrong end.

\textbf{16.} 2,2-dimethylpentane, 2,3-dimethylpentane,
2,4-dimethylpentane and 3,3-dimethylpentane.

\textbf{17.} Yes, 3-ethylpentane. On carbon 2 the ethyl group would make
a six-carbon chain through it: the \omterm{def:g7:atoms-and-molecules:molecule}{molecule} would be 3-methylhexane.

\textbf{18.} 2,2,3-trimethylbutane,
\chemfig{-[:0](-[:90])(-[:-90])-[:0](-[:90])-[:0]}.

\textbf{19.} $1 + 2 + 4 + 1 + 1 = 9$.

\textbf{20.} Heptane has 9 \omterm{def:g11:organic-skeletons:isomer}{constitutional isomers}, itself included.
\end{solution}
